% Optional information paths among four research objects.
\begin{tikzpicture}[
  x=1mm,y=1mm,font=\figurefont\fontsize{8.5}{11}\selectfont,
  box/.style={draw=BookInk,line width=.8pt,fill=BookPaper,
    minimum width=32mm,minimum height=12mm,align=center,inner sep=1.5mm},
  flow/.style={-{Latex[length=2mm]},draw=BookInk,line width=.9pt,
    rounded corners=1.2mm},
  model/.style={flow,dashed,draw=BookAmber},
  tag/.style={fill=white,inner sep=.7mm}
]
  \node[box] (data) at (54,61) {真实经验\\$\mathcal D$};
  \node[box] (value) at (26,30) {价值估计\\$\widehat V,\widehat Q$};
  \node[box,fill=BookAmber!8] (model) at (82,30) {环境模型\\$\widehat P,\widehat R$};
  \node[box,fill=BookTeal!10] (policy) at (54,-2) {策略\\$\pi$};
  \draw[flow] (data.south west) -- node[tag,above left] {预测监督} (value.north);
  \draw[flow] (data.south east) -- node[tag,above right] {拟合后果} (model.north);
  \draw[flow] (value.south) -- node[tag,below left] {策略更新} (policy.north west);
  \draw[model] (model.south) -- node[tag,below right] {模型内规划} (policy.north east);
  \draw[model] (model) -- node[tag,above] {模拟经验} (value);
  \draw[flow] (policy.south) -- (54,-14) -- (2,-14) -- (2,61) -- (data.west);
  \node[rotate=90,fill=white,inner sep=1mm] at (2,24) {真实交互后收集新经验};
\end{tikzpicture}
